% ----------------------------------------------------------
% Conclusão
% ----------------------------------------------------------
\chapter{Conclusão}

No último capítulo são discutidas opiniões dos integrantes acerca do que se tinha como expectativa e como de fato ocorreu, bem como observações do projeto no geral.

\section{Considerações Finais}
No escopo inicial, o contexto do projeto focou nas demandas de um pequeno negócio destacando a importância de um acervo digital e de ferramentas para a gestão interna. A fase de análise do sistema detalhou os requisitos de uma aplicação web voltada para o gerenciamento de inventário e de clientes, além da geração de relatórios de estoque e vendas.
Em seguida, o projeto do sistema estruturou os principais componentes da aplicação — envolvendo clientes, livros, histórico, estoque e sistema de sugestões. Essa modelagem foi traduzida em uma arquitetura tecnológica composta por React no frontend, Python no backend e SQLite para persistência.
Por fim, a etapa de implementação abordou as restrições técnicas e legais exigidas, como a conformidade com a LGPD, além de descrever o ambiente de desenvolvimento local via e a organização modular dos componentes. Observou-se também que o modelo conceitual passou por ajustes práticos durante a codificação, não se limitando de forma engessada aos diagramas iniciais. Dessa forma, os resultados confirmaram a viabilidade e a eficácia da solução desenvolvida, cumprindo as metas estabelecidas.

\section{Dificuldades Encontradas}
O tipo de aplicação abordada no projeto não se qualifica exatamente como um tipo de programa muito complexo. Não foram sentidas muitas complicações, no entanto o que se destaca na pequena quantidade observada, é que houveram alguns problemas de comunicação quanto ao que exatamente seria feito, uma dificuldade em encontrar uma organização com potencial para o desenvolvimento de um projeto interessante e dentro dos requisitos da matéria tanto quanto para escolher quais recursos usar em cada parte.

\section{Trabalhos Futuros}
Sistemas digitais podem possuir várias versões, assim como os diversos aplicativos que usamos em no dia-a-dia. Futuros incrementos para este software podem variar e acoplar todas as diferentes camadas de arquitetura. Em relação a funcionalidade, o sistema pode possuir menus tanto para funcionários quanto clientes que desejem visitar o sebo, pode-se implementar uma biblioteca mais adequada a partir do React, para gadgets e estilização ou mesmo adequar a solução para tornar possível escalá-la apropriadamente.